\documentclass{article}
\usepackage{amsmath,amssymb,amsthm}
\usepackage{algorithm}
\usepackage{algpseudocode}
\usepackage{booktabs}
\usepackage{hyperref}
\usepackage{geometry}
\geometry{margin=1in}
\newtheorem{theorem}{Theorem}
\newtheorem{lemma}{Lemma}
\newtheorem{definition}{Definition}
\newcommand{\prox}{\operatorname{prox}}
\newcommand{\grad}{\nabla}
\newcommand{\R}{\mathbb{R}}

\begin{document}

\section*{Proximal Gradient Method}

\section*{Mathematical Formulation}
We consider the composite convex optimization problem  
\[
\min_{x\in\R^{d}} \; F(x)\;:=\; f(x)+g(x),
\]
where  

\begin{itemize}
  \item $f:\R^{d}\!\to\!\R$ is convex and $L$‑smooth, i.e. $\grad f$ is $L$‑Lipschitz:
  \[
  \|\grad f(x)-\grad f(y)\|\le L\|x-y\|\qquad\forall x,y\in\R^{d}.
  \]
  \item $g:\R^{d}\!\to\!\R\cup\{+\infty\}$ is proper, closed and convex (possibly nonsmooth).
\end{itemize}
For a stepsize $\eta>0$, the \textbf{proximal gradient update} is  
\[
\boxed{%
x_{k+1}= \prox_{\eta g}\!\bigl(x_{k}-\eta\,\grad f(x_{k})\bigr)},
\qquad k=0,1,2,\dots
\tag{1}
\]
The proximal operator of $g$ is defined by  
\[
\prox_{\eta g}(u):=\arg\min_{z\in\R^{d}}\Bigl\{ g(z)+\frac{1}{2\eta}\|z-u\|^{2}\Bigr\}.
\tag{2}
\]

\section*{Pseudocode}
\begin{algorithm}[H]
\caption{Proximal Gradient Method}
\begin{algorithmic}[1]
\Require Initial point $x_{0}\in\R^{d}$, stepsize $\eta\in(0,1/L]$.
\For{$k=0,1,2,\dots$}
    \State Compute gradient $g_{k}\gets\grad f(x_{k})$.
    \State Gradient step $y_{k}\gets x_{k}-\eta g_{k}$.
    \State Proximal step $x_{k+1}\gets\prox_{\eta g}(y_{k})$.
\EndFor
\Ensure Final iterate $x_{K}$ (or any stopping criterion).
\end{algorithmic}
\end{algorithm}

\section*{Dry Run Example}
Minimize $f(w)=(w-3)^{2}$ (hence $g\equiv0$) with $\eta=0.1$ and $w_{0}=0$.

\begin{center}
\begin{tabular}{@{}c c c c c@{}}
\toprule
Iteration $k$ & $w_{k}$ & $\grad f(w_{k})=2(w_{k}-3)$ & $w_{k+1}=w_{k}-\eta\grad f(w_{k})$ & $f(w_{k+1})$\\
\midrule
0 & $0$ & $-6$ & $0-0.1(-6)=0.6$ & $(0.6-3)^{2}=5.76$\\
1 & $0.6$ & $-4.8$ & $0.6-0.1(-4.8)=1.08$ & $(1.08-3)^{2}=3.6864$\\
2 & $1.08$ & $-3.84$ & $1.08-0.1(-3.84)=1.464$ & $(1.464-3)^{2}=2.3665$\\
\bottomrule
\end{tabular}
\end{center}
Because $g=0$, $\prox_{\eta g}$ is the identity, so the update reduces to ordinary gradient descent.

\newpage
\section*{Assumptions}
\begin{enumerate}
  \item \textbf{Convexity of $f$ and $g$:} $\forall x,y,\;\lambda\in[0,1]$,
        \[
        f(\lambda x+(1-\lambda)y)\le \lambda f(x)+(1-\lambda)f(y),\qquad
        g(\lambda x+(1-\lambda)y)\le \lambda g(x)+(1-\lambda)g(y).
        \]
  \item \textbf{$L$‑smoothness of $f$:}
        \[
        f(y)\le f(x)+\langle\grad f(x),y-x\rangle+\frac{L}{2}\|y-x\|^{2},
        \qquad\forall x,y\in\R^{d}. \tag{A2}
        \]
  \item \textbf{Stepsize bound:} $\displaystyle 0<\eta\le\frac{1}{L}$.
  \item \textbf{Existence of a minimizer:} $F$ attains its minimum at some $x^{\star}\in\R^{d}$.
\end{enumerate}

\section*{Main Convergence Theorem}
\begin{theorem}[O(1/k) convergence of the proximal gradient method]\label{thm:main}
Let $\{x_{k}\}_{k\ge0}$ be generated by \eqref{1} with a stepsize $\eta\in(0,1/L]$ and let $x^{\star}$ be any minimizer of $F$. Then for all $K\ge1$,
\[
F(x_{K})-F(x^{\star})\;\le\;\frac{\|x_{0}-x^{\star}\|^{2}}{2\eta K}.
\tag{3}
\]
Equivalently, for the averaged iterate $\bar{x}_{K}:=\frac{1}{K}\sum_{k=1}^{K}x_{k}$,
\[
F(\bar{x}_{K})-F(x^{\star})\;\le\;\frac{\|x_{0}-x^{\star}\|^{2}}{2\eta K}.
\tag{4}
\]
\end{theorem}

\section*{Theoretical Properties}
\begin{itemize}
  \item \textbf{Stability:} The method is \emph{non‑expansive} because the proximal operator is firmly non‑expansive; consequently the iterates remain bounded whenever a minimizer exists.
  \item \textbf{Stepsize sensitivity:} Any $\eta\le1/L$ guarantees the descent inequality (see proof). Larger stepsizes may violate the key inequality and can lead to divergence.
  \item \textbf{Comparison to baselines:} When $g\equiv0$, the method reduces to classical gradient descent with the same $O(1/k)$ rate. When $g$ is the indicator of a convex set, the method becomes projected gradient descent, inheriting the same guarantee.
\end{itemize}

\section*{Discussion}
The theorem shows that the proximal gradient method attains the optimal sublinear rate for the class of convex, $L$‑smooth $f$ plus convex $g$. The proof crucially relies on the non‑expansiveness of $\prox_{\eta g}$ (Lemma~\ref{lem:nonexp}) and on the smoothness inequality \eqref{A2}. Extensions to accelerated schemes (FISTA) improve the rate to $O(1/k^{2})$, while stochastic variants require additional variance control.

\newpage
\section*{Preliminaries (Definitions and Lemmas)}
\begin{definition}[Proximal operator]\label{def:prox}
For $\eta>0$, $\prox_{\eta g}(u)$ is defined by \eqref{2}.  
It satisfies the \emph{optimality condition}
\[
\frac{1}{\eta}\bigl(u-\prox_{\eta g}(u)\bigr)\in\partial g\!\bigl(\prox_{\eta g}(u)\bigr).
\tag{5}
\]
\end{definition}

\begin{lemma}[Non‑expansiveness of the proximal operator]\label{lem:nonexp}
For any $u,v\in\R^{d}$,
\[
\bigl\|\prox_{\eta g}(u)-\prox_{\eta g}(v)\bigr\|\le\|u-v\|.
\tag{6}
\]
\end{lemma}
\begin{proof}
The proximal operator is firmly non‑expansive; see e.g. \cite[Prop.~12.27]{Bauschke2011}. $\square$
\end{proof}

\begin{lemma}[Smoothness inequality]\label{lem:smooth}
If $\grad f$ is $L$‑Lipschitz, then for all $x,y$,
\[
f(y)\le f(x)+\langle\grad f(x),y-x\rangle+\frac{L}{2}\|y-x\|^{2}.
\tag{7}
\]
\end{lemma}
\begin{proof}
Standard consequence of the integral form of the gradient Lipschitz property. $\square$
\end{proof}

\newpage
\section*{Main Proof (Step‑by‑Step Derivation)}
Let $x^{\star}$ be a minimizer of $F$.  
Denote $y_{k}=x_{k}-\eta\grad f(x_{k})$ and $x_{k+1}=\prox_{\eta g}(y_{k})$ as in \eqref{1}.

\begin{enumerate}
\item \textbf{Optimality condition of the proximal step.}  
From \eqref{5} with $u=y_{k}$ we obtain a subgradient $s_{k+1}\in\partial g(x_{k+1})$ such that
\[
s_{k+1}= \frac{1}{\eta}\bigl(y_{k}-x_{k+1}\bigr)
        =\frac{1}{\eta}\bigl(x_{k}-x_{k+1}\bigr)-\grad f(x_{k}).
\tag{8}
\]
\item \textbf{Upper bound on $f(x_{k+1})$.}  
Apply the smoothness inequality \eqref{7} with $x=x_{k}$ and $y=x_{k+1}$:
\[
f(x_{k+1})\le f(x_{k})+\langle\grad f(x_{k}),x_{k+1}-x_{k}\rangle
               +\frac{L}{2}\|x_{k+1}-x_{k}\|^{2}.
\tag{9}
\]
\item \textbf{Upper bound on $g(x_{k+1})$ via subgradient inequality.}  
Since $s_{k+1}\in\partial g(x_{k+1})$, convexity of $g$ yields
\[
g(x_{k+1})\le g(x^{\star})+\langle s_{k+1},x_{k+1}-x^{\star}\rangle.
\tag{10}
\]
\item \textbf{Combine \eqref{9} and \eqref{10}.}  
Adding the two inequalities gives
\[
\begin{aligned}
F(x_{k+1})-F(x^{\star})
&\le \langle\grad f(x_{k}),x_{k+1}-x_{k}\rangle
   +\frac{L}{2}\|x_{k+1}-x_{k}\|^{2}
   +\langle s_{k+1},x_{k+1}-x^{\star}\rangle .
\end{aligned}
\tag{11}
\]
\item \textbf{Replace $s_{k+1}$ using \eqref{8}.}  
Substituting \eqref{8} into the last inner product:
\[
\begin{aligned}
\langle s_{k+1},x_{k+1}-x^{\star}\rangle
&= \Big\langle\frac{1}{\eta}(x_{k}-x_{k+1})-\grad f(x_{k}),\,x_{k+1}-x^{\star}\Big\rangle\\
&= \frac{1}{\eta}\langle x_{k}-x_{k+1},x_{k+1}-x^{\star}\rangle
   -\langle\grad f(x_{k}),x_{k+1}-x^{\star}\rangle .
\end{aligned}
\tag{12}
\]
\item \textbf{Insert \eqref{12} into \eqref{11}.}  
The terms involving $\grad f(x_{k})$ cancel:
\[
\begin{aligned}
F(x_{k+1})-F(x^{\star})
&\le \frac{1}{\eta}\langle x_{k}-x_{k+1},x_{k+1}-x^{\star}\rangle
   +\frac{L}{2}\|x_{k+1}-x_{k}\|^{2}.
\end{aligned}
\tag{13}
\]
\item \textbf{Algebraic manipulation of the inner product.}  
Recall the identity $2\langle a,b\rangle =\|a\|^{2}+\|b\|^{2}-\|a-b\|^{2}$.  
With $a=x_{k}-x_{k+1}$ and $b=x_{k+1}-x^{\star}$ we obtain
\[
2\langle x_{k}-x_{k+1},x_{k+1}-x^{\star}\rangle
   =\|x_{k}-x^{\star}\|^{2}-\|x_{k+1}-x^{\star}\|^{2}
    -\|x_{k+1}-x_{k}\|^{2}.
\tag{14}
\]
Dividing \eqref{14} by $2\eta$ and inserting into \eqref{13} yields
\[
\begin{aligned}
F(x_{k+1})-F(x^{\star})
&\le \frac{1}{2\eta}\bigl(\|x_{k}-x^{\star}\|^{2}
          -\|x_{k+1}-x^{\star}\|^{2}
          -\|x_{k+1}-x_{k}\|^{2}\bigr)
   +\frac{L}{2}\|x_{k+1}-x_{k}\|^{2}.
\end{aligned}
\tag{15}
\]
\item \textbf{Use the stepsize condition $\eta\le 1/L$.}  
Since $L\eta\le1$, we have $ \frac{L}{2}\|x_{k+1}-x_{k}\|^{2}\le\frac{1}{2\eta}\|x_{k+1}-x_{k}\|^{2}$.  
Hence the two last terms in \eqref{15} combine to a non‑positive quantity:
\[
-\frac{1}{2\eta}\|x_{k+1}-x_{k}\|^{2}
   +\frac{L}{2}\|x_{k+1}-x_{k}\|^{2}
   \le 0.
\tag{16}
\]
\item \textbf{Descent inequality.}  
Discarding the non‑positive term (16) from \eqref{15} gives the key inequality
\[
\boxed{%
F(x_{k+1})-F(x^{\star})
\le \frac{1}{2\eta}\Bigl(\|x_{k}-x^{\star}\|^{2}
                         -\|x_{k+1}-x^{\star}\|^{2}\Bigr)
}
\tag{17}
\]
which is \textbf{[Step 1]} of the proof.
\item \textbf{Telescoping sum.}  
Sum \eqref{17} over $k=0,\dots,K-1$:
\[
\sum_{k=0}^{K-1}\!\bigl(F(x_{k+1})-F(x^{\star})\bigr)
\le \frac{1}{2\eta}\bigl(\|x_{0}-x^{\star}\|^{2}
                         -\|x_{K}-x^{\star}\|^{2}\bigr)
\le \frac{\|x_{0}-x^{\star}\|^{2}}{2\eta}.
\tag{18}
\]
\item \textbf{Extract the worst‑case bound.}  
Since the left‑hand side of \eqref{18} is a sum of $K$ non‑negative terms,
\[
K\bigl(F(x_{K})-F(x^{\star})\bigr)
\le \sum_{k=0}^{K-1}\!\bigl(F(x_{k+1})-F(x^{\star})\bigr)
\le \frac{\|x_{0}-x^{\star}\|^{2}}{2\eta}.
\tag{19}
\]
Dividing by $K$ yields exactly \eqref{3}, i.e. the $O(1/K)$ rate.
\item \textbf{Averaged iterate.}  
Convexity of $F$ implies
\[
F(\bar{x}_{K})\le\frac{1}{K}\sum_{k=1}^{K}F(x_{k}),
\]
and the same bound as in \eqref{19} follows, establishing \eqref{4}.
\end{enumerate}
Thus Theorem~\ref{thm:main} is proved. $\square$

\section*{Rate Derivation (With Constants)}
From \eqref{19} we have the explicit constant
\[
F(x_{K})-F(x^{\star})\le
\frac{\|x_{0}-x^{\star}\|^{2}}{2\eta K},
\qquad
\text{with }0<\eta\le\frac{1}{L}.
\]
Choosing the largest admissible stepsize $\eta=1/L$ minimizes the bound:
\[
F(x_{K})-F(x^{\star})\le\frac{L\|x_{0}-x^{\star}\|^{2}}{2K}.
\tag{20}
\]

\section*{Special Cases}
\begin{enumerate}
  \item \textbf{Pure gradient descent ($g\equiv0$).}  
        $\prox_{\eta g}$ reduces to the identity, and the proof collapses to the classical $O(1/k)$ result for smooth convex $f$.
  \item \textbf{Projected gradient descent.}  
        If $g$ is the indicator of a closed convex set $\mathcal{C}$,
        $\prox_{\eta g}= \Pi_{\mathcal{C}}$ (projection). The same bound \eqref{3} holds for the projected iterates.
  \item \textbf{Strongly convex $F$.}  
        If $F$ is $\mu$‑strongly convex, inequality \eqref{17} can be strengthened to a linear rate:
        \[
        F(x_{k+1})-F(x^{\star})\le\Bigl(1-\eta\mu\Bigr)\bigl(F(x_{k})-F(x^{\star})\bigr).
        \]
        Hence $F(x_{k})-F(x^{\star})\le(1-\eta\mu)^{k}\bigl(F(x_{0})-F(x^{\star})\bigr)$.
\end{enumerate}

\end{document}